\documentclass[11pt,reqno]{amsart}
\usepackage[T1]{fontenc}
\usepackage[utf8]{inputenc}
\usepackage{amsmath,amssymb,amsthm,mathtools}
\usepackage[margin=1.15in]{geometry}
\usepackage{xcolor}
\usepackage{booktabs}
\usepackage[expansion=false]{microtype}
\usepackage[colorlinks=true,linkcolor=blue!55!black,citecolor=blue!55!black,urlcolor=blue!55!black]{hyperref}
\usepackage{url}

% Typesetting only: lets TeX resolve tight lines by stretching interword
% space in a final pass. No effect on content.
\emergencystretch=2em

\theoremstyle{plain}
\newtheorem{theorem}{Theorem}[section]
\newtheorem{proposition}[theorem]{Proposition}
\newtheorem{lemma}[theorem]{Lemma}
\newtheorem{corollary}[theorem]{Corollary}
\theoremstyle{definition}
\newtheorem{definition}[theorem]{Definition}
\newtheorem{remark}[theorem]{Remark}
\newtheorem{example}[theorem]{Example}
\newtheorem{question}[theorem]{Question}
\newtheorem*{theoremR}{Theorem R}
\newtheorem*{theoremS}{Theorem S}

\DeclareMathOperator{\lcp}{lcp}
\DeclareMathOperator{\ice}{ice}
\DeclareMathOperator{\Dio}{Dio}
\newcommand{\N}{\mathbb{N}}
\newcommand{\Z}{\mathbb{Z}}
\newcommand{\Q}{\mathbb{Q}}
\newcommand{\R}{\mathbb{R}}
\newcommand{\Zt}{\Z_2}
\newcommand{\Qt}{\Q_2}
\newcommand{\PH}{\Phi}
\newcommand{\gs}{\gamma^{*}}
\newcommand{\Av}{A(\gamma)}
\newcommand{\om}{\omega_1^{(2)}}
\newcommand{\vtwo}{v_2}

\begin{document}

\title[Transcendence of the $3x+1$ conjugacy map on Sturmian words]
      {Transcendence of the $3x+1$ conjugacy map on Sturmian words}
\author{Elias De Jes\'us}
\address{Independent Researcher}
\email{dejesuselias10@gmail.com}
\urladdr{https://orcid.org/0009-0007-0190-9143}
\date{October 2026}

\subjclass[2020]{Primary 11J61; Secondary 11B83, 68R15, 11J81}
\keywords{$3x+1$ problem, conjugacy map, Sturmian word, $2$-adic transcendence,
Ridout's theorem, initial critical exponent}

\begin{abstract}
Let $T$ be the $3x+1$ map on the $2$-adic integers and let $\PH\colon\Zt\to\Zt$ be the
Bernstein--Lagarias conjugacy map, which sends a parity vector to the unique $2$-adic integer
realising it. We prove that $\PH(c_\gamma)$, $\PH(1c_\gamma)$, $\PH(0c_\gamma)$ and $\PH$ of every
shift of these words are transcendental for every irrational slope
$\gamma<\gs=(3+\sqrt5)/(4\log_2 3)=0.8258977\ldots$, with no hypothesis whatever on the partial
quotients of $\gamma$; in particular at the resonance slope $\gamma=\log_3 2$. The mechanism is a
bridge: by the Bernstein--Lagarias isometry, the $2$-adic approximation depth of a periodic shadow
equals a longest common prefix, so the Diophantine behaviour of $\PH$ is governed by the repetition
structure of the parity word, at a height cost $\Av=\max(1,\gamma\log_2 3)$ that we control at every
level. Quantitatively $\om(\PH(c_\gamma))\ge\ice(c_\gamma)/\Av-1$, where $\ice$ is the initial
critical exponent of Berth\'e--Holton--Zamboni; their floor $\ice(c_\gamma)\ge1+\varphi$ then beats
the Ridout threshold $\om\le1$ by a margin bounded away from zero. The result is not a corollary of
the digit-based transcendence criteria: the $2$-adic Hensel digits of $\PH(c_{\log_3 2})$ are
measured to have full complexity $2^n$, against the parity word's $n+1$, and the shadow denominators
$2^{\ell}-3^{k}$ are not $S$-units for any fixed finite $S$. We also record the Liouville
alternative for unbounded partial quotients and a finite effective irrationality measure for bounded
ones at slopes $\gamma\le\log_3 2$.
\end{abstract}

\maketitle

\section{Introduction}\label{sec:intro}

\subsection{The conjugacy map and the Periodicity Conjecture}
The $3x+1$ map is defined on the $2$-adic integers $\Zt$ by
\[
T(x)=\frac{3x+1}{2}\quad(x\text{ odd}),\qquad T(x)=\frac{x}{2}\quad(x\text{ even}),
\]
and the $2$-adic shift by $S(x)=(x-1)/2$ for $x$ odd and $S(x)=x/2$ for $x$ even, so that $S$ shifts
the $2$-adic digit expansion. Bernstein and Lagarias \cite[\S1]{BL96} proved that $T$ and $S$ are
conjugate: there is a unique homeomorphism $\PH\colon\Zt\to\Zt$ with $\PH(0)=0$ and
\begin{equation}\label{eq:conj}
  \PH\circ S\circ\PH^{-1}=T .
\end{equation}
Identifying a word $v=v_0v_1v_2\cdots\in\{0,1\}^{\N}$ with the $2$-adic integer $\sum_k v_k2^k$, the
inverse $\PH^{-1}(x)=\sum_{k\ge0}\bigl(T^k(x)\bmod 2\bigr)2^k$ is the \emph{parity vector} of $x$,
so $\PH(v)$ is the unique $x\in\Zt$ whose parity vector is $v$, and \eqref{eq:conj} reads
\begin{equation}\label{eq:shift}
  \PH(\sigma v)=T(\PH(v)),
\end{equation}
$\sigma$ being the shift on words. A $2$-adic integer is rational if and only if its digit
expansion is eventually periodic; and if $v$ is eventually periodic then $\PH(v)\in\Q$, by
Proposition~\ref{prop:shadow} below. The \emph{Periodicity Conjecture}, which goes back to Lagarias's survey \cite{Lagarias85} and is
stated in the form below in \cite[Abstract]{LS09} and \cite[\S1]{LS21}, asserts the converse:

\begin{question}[Periodicity Conjecture; open]\label{q:per}
Is the parity vector of every rational $2$-adic integer eventually periodic? Equivalently: is
$\PH(v)$ irrational for every aperiodic $v\in\{0,1\}^{\N}$?
\end{question}

Nothing in the present paper bears on Question~\ref{q:per} in general: it quantifies over an
uncountable family of words. What can be done is to settle it, and more, on thin but genuinely
aperiodic families. The thinnest are the Sturmian words, of factor complexity $p(L)=L+1$.

\subsection{Prior irrationality results}
For Sturmian words Question~\ref{q:per} is a theorem. Write $\beta=\log_3 2=0.6309297\ldots$ for the
\emph{resonance} slope. Three strands of prior work:

\begin{itemize}
\item Monks and Yazinski \cite[Theorem 2.7(b)]{MY04} proved that a rational $2$-adic integer with
  infinite orbit has lower ones-density at least $\beta$. Contrapositively, $\PH(v)\notin\Q$ for
  every aperiodic $v$ of lower ones-density $<\beta$; in particular for every Sturmian word of slope
  $\gamma<\beta$.
\item L\'opez and Stoll \cite[\S4]{LS09} proved that the \emph{real} value
  $\PH_{\R}(m_\alpha)$ is irrational for $1>\alpha>\beta$, using the continued fraction of their
  Corollary~2, with a dual construction below $\beta$. This is a different number from the $2$-adic
  $\PH$; see Remark~\ref{rem:real}. Their later preprint \cite[Theorem 1]{LS21} concerns the
  $2$-adic value; we cite it as context only and depend on it nowhere.
\item The author \cite[Theorems 8.5 and 10.1, Corollary 10.5]{DJirr} proved that $\PH(s)\notin\Q$
  for \emph{every} mechanical word $s$ of every irrational slope and every intercept, effectively.
  Cassidy \cite[Theorems R and S]{Cassidy26} proved the irrationality statement independently and
  contemporaneously, by a $2$-adic repetition criterion, for every slope and intercept and for a
  wider class of maps; his note was posted first (see \S\ref{sec:related}).
\end{itemize}

Irrationality says $\PH(s)$ is not a rational point. It says nothing about whether $\PH(s)$ is
algebraic. That is the gap this paper closes.

\subsection{Main results}
Let $\gamma\in(0,1)$ be irrational with continued fraction $\gamma=[0;a_1,a_2,\dots]$ and
convergent denominators $q_k$. Let $c_\gamma$ be the characteristic Sturmian word of slope $\gamma$
and let $1c_\gamma$, $0c_\gamma$ be the two mechanical words of slope $\gamma$ at intercept $0$.
Put
\[
  \Av:=\max(1,\gamma\log_2 3),\qquad
  \gs:=\frac{1+\varphi}{2\log_2 3}=\frac{3+\sqrt5}{4\log_2 3}=0.8258977696818354644\ldots,
\]
where $\varphi=(1+\sqrt5)/2$, and let $\ice$ denote the initial critical exponent of
Berth\'e--Holton--Zamboni \cite[\S2]{BHZ06} (Definition~\ref{def:ice}). Our normalisation of the
$2$-adic Koksma exponent $\om$ is fixed in \S\ref{ss:norm}.

\begin{theorem}[Quantitative form]\label{thm:quant}
For every irrational $\gamma\in(0,1)$,
\[
  \om\bigl(\PH(c_\gamma)\bigr)\;\ge\;\frac{\ice(c_\gamma)}{\Av}-1
  \;=\;\frac{1+\limsup_{k}q_{k+1}/q_k}{\Av}-1 .
\]
\end{theorem}

\begin{theorem}[Transcendence]\label{thm:main}
Let $\gamma\in(0,1)$ be irrational and suppose
\begin{equation}\label{eq:T}
  \ice(c_\gamma)>2\Av . \tag{T}
\end{equation}
Then $\PH(c_\gamma)$ is transcendental, and so are $\PH(1c_\gamma)$, $\PH(0c_\gamma)$ and
$\PH(\sigma^k 1c_\gamma)$ for every $k\ge0$.
\end{theorem}

\begin{corollary}[Unconditional range]\label{cor:uncond}
For every irrational $\gamma<\gs$ the numbers $\PH(c_\gamma)$, $\PH(1c_\gamma)$, $\PH(0c_\gamma)$
and $\PH$ of every shift of these words are transcendental. In particular they are transcendental at
the resonance slope $\gamma=\beta=\log_3 2$, \emph{with no hypothesis on the partial quotients of
$\log_2 3$}.
\end{corollary}

Two complements. If the partial quotients of $\gamma$ are unbounded then $\ice(c_\gamma)=\infty$ and
the conclusion sharpens to a Liouville statement (Theorem~\ref{thm:liouville}); if they are bounded
and $\gamma\le\beta$ then $\PH(c_\gamma)$ has a finite \emph{effective} irrationality measure
(Theorem~\ref{thm:Bstar}).

Theorem~\ref{thm:main} and Corollary~\ref{cor:uncond} answer problem~(3) of \cite[\S12]{DJirr}
--- ``is $\PH(1c_\beta)$, or $\PH(1c_\gamma)$ for a given irrational $\gamma$, transcendental
over $\Q$?'' --- for every irrational $\gamma<\gs$, the resonance slope $\gamma=\beta$ included.
Problem~(2) of \cite[\S12]{DJirr} is a different question and is \emph{not} answered here: it asks
for $\mu<\infty$ and $C>0$ with $|\PH(1c_\gamma)-y|_2\ge C|y|^{-\mu}$ for all \emph{integers}
$y\ge1$. Theorems~\ref{thm:liouville} and \ref{thm:Bstar} concern approximation by
\emph{rationals} and do not bear on it; \S\ref{ss:integer} records only the negative fact that the
Liouville blow-up of Theorem~\ref{thm:liouville} does not descend to $\Z$, so problem~(2) is not
obstructed by it. It remains open.

\subsection{Non-claims}\label{ss:nonclaims}
Stated plainly, so that the scope is not inferred from the title. This paper does \emph{not} bear on
the convergence of the $3x+1$ map on the positive integers, nor on the existence of divergent or
cyclic orbits. It does not prove the Periodicity Conjecture (Question~\ref{q:per}), which quantifies
over an uncountable family of aperiodic parity vectors; the Sturmian words form a measure-zero,
countable-slope subfamily. It proves nothing for slopes $\gamma\ge\gs$ except what
Theorem~\ref{thm:main} gives under its hypothesis~\eqref{eq:T}, and $\gs$ is sufficient, not
necessary: slopes above $\gs$ with large partial quotients still satisfy \eqref{eq:T}
(\S\ref{ss:aboveg}). It says nothing about intercepts other than $0$ beyond \S\ref{ss:intercepts},
and nothing about approximation by \emph{integers}, which is problem~(2) of \cite[\S12]{DJirr} and
remains open (\S\ref{ss:integer}). No claim of priority is made over \cite{Cassidy26}; see
\S\ref{sec:related}.

\subsection{Epistemic tiers}\label{ss:tiers}
Three labels are used, and used consistently. \textbf{T1}: proved here or quoted from a cited source
with its hypotheses checked --- Theorems~\ref{thm:quant}--\ref{thm:main},
Corollary~\ref{cor:uncond}, Propositions~\ref{prop:iso}--\ref{prop:irr},
Theorems~\ref{thm:liouville} and \ref{thm:Bstar}, and the two propositions of
\S\ref{sec:truncations}. \textbf{T2}: established by exact integer or $2$-adic computation, with the
precision stated --- every table and every numerical assertion in Appendix~\ref{app:verif},
Example~\ref{ex:witness}, the complexity table of \S\ref{ss:notcor}, and the quartic witness in the
footnote of \S\ref{ss:norm}. No floating-point number decides any inequality anywhere in this paper.
\textbf{T3}: interpretation, flagged as such in the text --- the reading of \eqref{eq:terms} as the
reason the archimedean place is unavailable, and the discussion in \S\ref{sec:scope}.

\subsection{Why this is not a corollary of the digit-based criteria}\label{ss:notcor}
There is a well-developed theory which concludes transcendence from low complexity of a \emph{digit
expansion}: Ferenczi--Mauduit \cite{FM97} for low-complexity $b$-ary expansions, the combinatorial
transcendence criterion of Adamczewski--Bugeaud \cite{AB07} and its $p$-adic analogues, and --- the
closest formal match --- Ooto \cite{Ooto16}, who proves that an irrational $p$-adic number whose
\emph{Hensel digits} form an automatic, primitive morphic or Sturmian sequence is an $S$-, $T$- or
$U_1$-number. Had any of these applied, Theorem~\ref{thm:main} would be a corollary, and Ooto's
theorem would in addition supply the Mahler class. None applies, for two independent reasons.

\emph{First, the wrong coordinate.} The Sturmian structure of our object lives in the parity vector,
not in the digits, and $\PH$ destroys it on the way to the digits. This is measurable. For
$\gamma=\log_3 2$ we computed the first $3000$ Hensel digits of $\PH(c_\gamma)$ in exact arithmetic
from the parity-vector definition and counted distinct factors:

\begin{center}
\begin{tabular}{lcccccl}
\toprule
factor length $n$ & $2$ & $3$ & $4$ & $5$ & $6$ \\
\midrule
Hensel digits of $\PH(c_{\log_3 2})$ & $4$ & $8$ & $16$ & $32$ & $64$ \\
the parity word $c_{\log_3 2}$       & $3$ & $4$ & $5$  & $6$  & $7$  \\
Sturmian requires                    & $3$ & $4$ & $5$  & $6$  & $7$  \\
\bottomrule
\end{tabular}
\end{center}

\noindent
The digits attain the maximum $2^n$ at every length computed, while the parity word is Sturmian on
the nose. So the hypothesis of each criterion above fails as badly as it can. (This is also consistent with
the closing conjecture of \cite{LS09}, that $\PH$ "always maps Sturmian words to infinite words of
full complexity", at lengths $\le6$ for this slope.)

\emph{Second, the denominators are not $S$-units.} Our approximants are the periodic shadows
$\PH(W^\infty)=c_W/(2^{\ell}-3^{k})$ of Proposition~\ref{prop:shadow}. No fixed finite set $S$ of
primes contains the prime divisors of $2^{\ell}-3^{k}$ for all $\ell,k$. In the Adamczewski--Bugeaud
construction the denominators are $b^{s}-b^{r}$, supported on the primes of $b$ together with
$b^{s-r}-1$, and Ridout's $S$-unit refinement buys a full unit of exponent when the denominator is
an $S$-unit. We get no such gain. This is exactly why the threshold in \eqref{eq:T} is $2\Av$ and
not $\Av$, and why initial squares --- $\ice\ge2$, which holds for every Sturmian word
\cite[Theorem 1]{ADQZ2001} --- are \emph{not} enough at $\Av=1$: they give $\om\ge1$, the Ridout
borderline, which proves nothing. The entire margin of the proof is the amount by which
Berth\'e--Holton--Zamboni's floor $1+\varphi=2.6180\ldots$ exceeds $2$.

\subsection{The bridge}\label{ss:bridge}
The contribution of this paper, and of \cite{DJirr} before it, is best described as a bridge rather
than a technique. The Bernstein--Lagarias isometry (Proposition~\ref{prop:iso}) says
\[
  \vtwo\bigl(\PH(v)-\PH(w)\bigr)=\lcp(v,w),
\]
so the $2$-adic approximation depth of a periodic shadow is \emph{exactly} a longest common prefix.
A repetition in the parity word is therefore a Diophantine approximation, with no loss. What the
bridge costs is height: the shadow of a block of length $\ell$ has
$\log_2 H\le\Av\,\ell+O(\log\ell)$ (Proposition~\ref{prop:height}). Hence
\[
  \text{Diophantine exponent of }\PH(s)\;\longleftrightarrow\;
  \frac{\text{repetition exponent of }s}{\Av},
\]
and any theorem about initial repetitions in Sturmian words becomes a theorem about $\PH$. The two
payoffs recorded so far are irrationality \cite{DJirr}, which needs only that the repetitions are
unbounded, and transcendence (this paper), which needs them to beat $2\Av$.

The two differ by exactly one factor of two in the exponent, and that factor is the whole of the
distance between the two papers. Irrationality \cite[\S10]{DJirr} needs the initial repetition
exponent of $s$ to exceed $\Av$: a shadow of height $2^{\Av\ell+O(\log\ell)}$ agreeing with
$\PH(s)$ to depth $>\Av\ell$ is already incompatible with a fixed rational value. Transcendence
needs the exponent to exceed $2\Av$, because Theorem R --- the $2$-adic Ridout theorem --- bites
only beyond approximation exponent $2$, equivalently $\om>1$ in the normalisation of
\S\ref{ss:norm}. What supplies the second, and with room to spare, is the
Berth\'e--Holton--Zamboni floor $\ice(c_\gamma)\ge1+\varphi=2.6180\ldots$
(Proposition~\ref{prop:icefloor}); the inequality $1+\varphi>2\Av$ is exactly the condition
$\gamma<\gs$ of Corollary~\ref{cor:uncond}.

\subsection{What is new here}
New in this paper: the transcendence theorem and its range, Theorem~\ref{thm:main} and
Corollary~\ref{cor:uncond}; the quantitative bound $\om\ge\ice/\Av-1$ of
Theorem~\ref{thm:quant}, obtained through $c_\gamma$; the parity obstruction of
\S\ref{sec:parity} with the witness of Example~\ref{ex:witness}; and the two complements,
Theorem~\ref{thm:liouville} and Theorem~\ref{thm:Bstar}.

Three ingredients are the author's \emph{prior} work. They are quoted here with exact locators and
are not claimed as new: the depth law at both convergent parities is \cite[Theorem 8.3]{DJirr}
(used in \S\ref{sec:parity}, and see Remark~\ref{rem:depthfromP1}); the height bound
$c_W\le3\ell\max(2^{\ell},3^{k})$ for cyclic permutations of standard words is
\cite[Lemma 10.4]{DJirr} (Proposition~\ref{prop:height}); and the transfer identities are
\cite[Corollaries 7.1 and 8.8]{DJirr} (equation~\eqref{eq:affine}). Everything else is inherited:
in particular the isometry is \cite{BL96}, the periodic shadows and their approximation depths are
\cite[Theorem 1]{LS09}, the term-size estimate is \cite[Lemma 43, \S11]{LS21}, the $\ice$ formula
and its floor are \cite[\S4.2]{BHZ06}, and the threshold is \cite[Theorem 1.3]{BK18}. Where we
keep a proof of an inherited statement it is labelled a re-derivation. A table of every ingredient
against its prior source accompanies this paper.

Transcendence of $\PH(s)$ for every Sturmian word $s$ was first proposed by Cassidy
\cite[\S4.6]{Cassidy26}, with a sketch via Schlickewei's $p$-adic Subspace Theorem, recorded
there as unaudited. We give a complete proof for characteristic words of every irrational slope
$\gamma<\gs$, by a different route: Ridout's theorem in the single-prime form
\cite[Theorem 1.3]{BK18}, together with the Berth\'e--Holton--Zamboni bound on initial powers.
The general claim remains open.


\subsection{Organisation}
\S\ref{sec:setup} fixes notation and assembles the inherited apparatus. \S\ref{sec:main} proves
Theorems~\ref{thm:quant} and \ref{thm:main}. \S\ref{sec:parity} explains, with an explicit witness,
why the argument must run through $c_\gamma$. \S\ref{sec:liouville} records the Liouville
alternative, \S\ref{sec:measure} a finite effective measure in the bounded-partial-quotient case.
\S\ref{sec:scope} lists the scope limitations and open problems, \S\ref{sec:related} the related
work. Appendix~\ref{app:schlickewei} re-derives the threshold from Schlickewei's subspace theorem;
Appendix~\ref{app:verif} summarises the exact-arithmetic verification.

\section{Setup}\label{sec:setup}

\subsection{The conjugacy map and the isometry}
Throughout, $\vtwo$ is the $2$-adic valuation, $|x|_2=2^{-\vtwo(x)}$, and for words
$v,w\in\{0,1\}^{\N}$, $\lcp(v,w)$ is the length of their longest common prefix. For a word $v$ and
$m\ge0$ set $k_m(v):=\#\{i<m:v_i=1\}$.

\begin{proposition}[{\cite{BL96}; also recorded as \cite[eq.~(4), \S1]{LS21}}]\label{prop:iso}
For all $v,w\in\{0,1\}^{\N}$,
\[
  \vtwo\bigl(\PH(v)-\PH(w)\bigr)=\lcp(v,w),
  \qquad\text{equivalently}\qquad |\PH(x)-\PH(y)|_2=|x-y|_2 .
\]
In particular $\PH$ is injective: distinct words have finite $\lcp$, hence distinct values.
\end{proposition}

L\'opez and Stoll state the isometry as their equation~(4) and attribute it to \cite{BL96}; it is
restated as \cite[Proposition 2.2]{DJirr}. We also use, from \cite[\S1]{BL96}, that $\PH(v)$ is odd
if and only if $v_0=1$, and the identity \eqref{eq:shift}; iterating the latter and writing
$c_m(v):=\sum_{i<m,\,v_i=1}3^{\,k_m(v)-k_{i+1}(v)}2^{i}\in\Z$ gives
\begin{equation}\label{eq:affinegen}
  2^{m}\,\PH(\sigma^{m}v)=3^{\,k_m(v)}\,\PH(v)+c_m(v)\qquad(m\ge0),
\end{equation}
a relation with rational coefficients.

\begin{proposition}[{Periodic shadows; \cite[Theorem 1 and Lemma 12]{LS09}, restated as
\cite[Proposition 4.1]{DJirr}}]\label{prop:shadow}
Let $w$ be a finite word of length $\ell$ with $k\ge1$ ones. Then $2^{\ell}\ne3^{k}$ and
\[
  \PH(w^{\infty})=\frac{c_w}{2^{\ell}-3^{k}},\qquad
  c_w=\sum_{i<\ell,\,w_i=1}3^{\,k-k_{i+1}(w)}2^{i}\in\Z_{>0},
\]
and the denominator is \textbf{odd}.
\end{proposition}

The periodic approximants, and their approximation depths, go back to L\'opez and Stoll. Their
\cite[Theorem 1]{LS09} gives the closed $2$-adic series
\begin{equation}\label{eq:LS09series}
  \PH(1c_{\beta})=-\frac13-\sum_{j\ge0}(-1)^{j+1}
  \frac{2^{\,q_{j+1}+q_j-1}}{3\,(3^{p_{j+1}}-2^{q_{j+1}})(3^{p_j}-2^{q_j})},
\end{equation}
indexed by consecutive pairs of convergents of $\beta$, with \cite[Corollary 2]{LS09} the associated
generalised continued fraction for $-1/\PH(1c_\beta)$. The term exponents $q_{j+1}+q_j-1$ of
\eqref{eq:LS09series} \emph{are} the approximation depths used in this paper, at the convergents
from above; this was pointed out in \cite[\S1]{DJirr}. We verified \eqref{eq:LS09series} against an
independent computation of $\PH$ at three slopes, exactly modulo $2^{4000}$
(Appendix~\ref{app:verif}).

\subsection{Sturmian words, heights, and the constant \texorpdfstring{$\Av$}{A(gamma)}}
For irrational $\gamma\in(0,1)$ the characteristic word is $c_\gamma(j)=\lfloor(j+1)\gamma\rfloor
-\lfloor j\gamma\rfloor$ for $j\ge1$, and $1c_\gamma$, $0c_\gamma$ are $1$, resp.\ $0$, followed by
$c_\gamma$; these are the two mechanical words of slope $\gamma$ at intercept $0$. Every Sturmian
word is balanced, so any factor of length $\ell$ has $k$ ones with $|k-\gamma\ell|<1$. For
$x=u/v\in\Q$ in lowest terms with $v>0$ we write $H(x)=\max(|u|,v)$.

A finite word $W$ is a conjugate (cyclic permutation) of a standard word if and only if $W^\infty$
is balanced \cite[Ch.~2]{Lothaire02}. The input we need is:

\begin{proposition}[{\cite[Proposition 3.2]{BHZ06}}]\label{prop:bhz32}
If a Sturmian sequence begins in $W^{r}$ with $r\ge2$, $|W|\ge2$ and $W$ primitive, then $W$ is a
conjugate of a standard word; hence $W^{\infty}$ is balanced. If $W$ is not primitive, $W=V^{t}$ and
the sequence begins in $V^{rt}$, so primitivity costs nothing.
\end{proposition}

The height of a shadow is controlled by an archimedean estimate on the terms of the defining series.
At the resonance slope that estimate is \cite[Lemma 43, \S11]{LS21}: writing
$\lfloor j\alpha\rfloor=j\alpha-\{j\alpha\}$ with $\alpha=\log_2 3$ and using $2^{\alpha}=3$,
\begin{equation}\label{eq:terms}
  3^{-(j+1)}2^{\lfloor j\alpha\rfloor}=\tfrac13\,2^{-\{j\alpha\}}\in\bigl(\tfrac16,\tfrac13\bigr]
  \qquad (j\ge0),
\end{equation}
and L\'opez and Stoll give in addition the exact distribution of these terms inside
$(\frac16,\frac13]$, by an explicit permutation sorting them into
$\frac16<y_1<\cdots<y_{p_k}=\frac13$; see also \cite[Lemmas 41--42]{LS21} for their arithmetic and
geometric means. Read archimedean-ly, \eqref{eq:terms} is why the real series diverges at the
resonance, so that the $2$-adic place is the only one available; see \S\ref{sec:scope}. The
extension of the estimate from the resonance slope to every slope, by balance, is
\cite[Lemma 10.4]{DJirr}, the author's prior work, quoted here and not claimed as new:

\begin{proposition}[{Height of a shadow; \cite[Lemma 10.4]{DJirr}, resting on
\cite[Lemma 43]{LS21}}]\label{prop:height}
Let $W$ be a conjugate of a standard word, $|W|=\ell$, with $k\ge1$ ones, and suppose $W$ is a
factor of a Sturmian word of slope $\gamma$. Then $0<c_W\le3\ell\max(2^{\ell},3^{k})$ and
\[
  \log_2 H\bigl(\PH(W^{\infty})\bigr)\;<\;\Av\,\ell+\log_2(3\ell)+\log_2 3 .
\]
\end{proposition}

\begin{proof}
With $\delta=2^{\ell}-3^{k}$ and $g=\gcd(c_W,|\delta|)$ we have $H=\max(c_W,|\delta|)/g\le
\max(c_W,|\delta|)$. Since $2^{\ell}$ and $3^{k}$ are distinct positive integers,
$|\delta|<\max(2^{\ell},3^{k})$; and $c_W\le3\ell\max(2^{\ell},3^{k})$ by
\cite[Lemma 10.4]{DJirr}, whose input is that $W^{\infty}$ is balanced, so that
$|k_{i+1}(W)-(i+1)k/\ell|<1$. Hence $\log_2 H\le\log_2(3\ell)+\max(\ell,k\log_2 3)$. Balance of the
ambient Sturmian word gives $|k-\gamma\ell|<1$, so $k\log_2 3<\gamma\ell\log_2 3+\log_2 3$ and
therefore $\max(\ell,k\log_2 3)<\Av\ell+\log_2 3$.
\end{proof}

\subsection{The initial critical exponent}
\begin{definition}[{\cite[\S2]{BHZ06}}]\label{def:ice}
The \emph{prefix power} of a finite word $W$ in a sequence $\omega$ is the largest real $p$ such
that $W^{p}$ is a prefix of $\omega$; equivalently $\lcp(\omega,W^{\infty})/|W|$. The
\emph{initial critical exponent} $\ice(\omega)$ is the limit superior of the prefix powers of the
words $\omega[0,n)$ in $\omega$.
\end{definition}

Because the prefix power is defined as the largest \emph{real} such $p$, we have
$\lcp(\omega,W^{\infty})=p\,|W|$ \emph{exactly}: no floor and no $O(1)$ is lost in passing between
$\ice$ and approximation depth. This is the precise sense in which $\ice$ is the right quantity for
the bridge of \S\ref{ss:bridge}.

\begin{proposition}[{\cite[\S4.2]{BHZ06}}]\label{prop:icefla}
For the characteristic Sturmian sequence of slope $\alpha=[0;a_1,a_2,\dots]$,
\[
  \ice(c_\alpha)=\limsup_{k\to\infty}\frac{q_{k+1}+q_k-2}{q_k}
  =1+\limsup_{k\to\infty}\frac{q_{k+1}}{q_k}
  =1+\limsup_{k\to\infty}[a_{k+1};a_k,\dots,a_1].
\]
\end{proposition}

\begin{remark}[The $c_\gamma$ depth from {\cite[Theorem 8.3]{DJirr}}]\label{rem:depthfromP1}
The depth $q_{n+1}+q_n-2$ of Proposition~\ref{prop:icefla} can be read off the author's depth law
for $1c_\gamma$. Let $w_n$ be the periodic word attached to the convergent $p_n/q_n$ as in
\cite[\S8]{DJirr}. Then \cite[Theorem 8.3]{DJirr} gives
$\lcp(1c_\gamma,w_n^{\infty})=q_n+q_{n+1}-1$ for $n$ odd. Since $c_\gamma=\sigma(1c_\gamma)$
and the two words agree in their first letter, shifting removes exactly one position from the common
prefix, so
\[
  \lcp\bigl(c_\gamma,\sigma(w_n^{\infty})\bigr)=q_n+q_{n+1}-2
  \qquad(n\ \text{odd}),
\]
and $\sigma(w_n^{\infty})=(w_n')^{\infty}$ where $w_n'$ is the cyclic permutation of $w_n$ by one
position --- the same length, the same number of ones, hence a shadow of the same shape
(Proposition~\ref{prop:shadow}) with the same denominator bound. This recovers the depths of
Proposition~\ref{prop:icefla} along the odd levels, and with them
$\ice(c_\gamma)\ge1+\limsup_{n\ \mathrm{odd}}q_{n+1}/q_n$.

It does \emph{not} recover the even levels: there \cite[Theorem 8.3]{DJirr} gives
$\lcp(1c_\gamma,w_n^{\infty})=q_n-1$, so the shift yields only $q_n-2$. The limit superior over
\emph{all} $k$ in Proposition~\ref{prop:icefla} is Berth\'e--Holton--Zamboni's and is what we
quote; its approximant at every level $k$ is $W_k^{\infty}$ for $W_k$ the length-$q_k$ prefix of
$c_\gamma$ itself --- the word used in the proof of Theorem~\ref{thm:Bstar} --- and not the shift
$\sigma(w_n^{\infty})$ of a $1c_\gamma$ approximant. The asymmetry is not a defect of the
bookkeeping: it is the parity phenomenon of \S\ref{sec:parity}, and it is precisely what separates
$\ice(1c_\gamma)$ from $\ice(c_\gamma)$ there.
\end{remark}

\begin{proposition}[{Floor; \cite[\S4.2]{BHZ06}}]\label{prop:icefloor}
For every irrational $\gamma$,
\[
  \ice(c_\gamma)\;\ge\;1+\varphi=\varphi^{2}=\frac{3+\sqrt5}{2}=2.6180339887\ldots,
\]
with equality if and only if $a_k=1$ for all large $k$.
\end{proposition}

This is \cite[\S4.2]{BHZ06}: in the sentence following the proof of their Theorem~1.2 they state
that $\ice(\omega)\le3$ if and only if all but finitely many $a_k$ equal $1$, in which case
$\ice(\omega)=1+\theta$, where $\theta=(1+\sqrt5)/2$ is the golden mean fixed on \cite[p.~3]{BHZ06}.
Eventually-all-ones therefore gives $\ice=1+\varphi$; otherwise $\ice>3>1+\varphi$. The Fibonacci
equality case is also stated in their introduction (the characteristic sequence ``begins in no
$(3+\sqrt5)/2\approx2.62$ power at all'', so the supremum is not attained), and the constant
$1+\theta$ appears in \cite[Proposition 2.1(3)]{BHZ06}. The formula of
Proposition~\ref{prop:icefla} is also obtainable from Cassaigne's recurrence quotient
\cite{Cassaigne99}; the companion index formula is Vandeth \cite[Theorem 16]{Vandeth00} and the
Fibonacci repetition constants are Mignosi--Pirillo \cite{MignosiPirillo92}. For completeness we
record the elementary re-derivation of the floor; it is not a new result.

\begin{lemma}[Re-derivation of Proposition~\ref{prop:icefloor}]\label{lem:phifloor}
For every irrational $\gamma=[0;a_1,a_2,\dots]$ one has $\limsup_k q_{k+1}/q_k\ge\varphi$, with
equality if and only if $a_k=1$ for all large $k$.
\end{lemma}

\begin{proof}
$q_{k+1}/q_k=a_{k+1}+q_{k-1}/q_k>a_{k+1}$. Suppose $L:=\limsup_k q_{k+1}/q_k<\varphi$. As
$\varphi<2$, for all large $k$ we get $a_{k+1}<2$, i.e.\ $a_{k+1}=1$; then $q_{k+1}=q_k+q_{k-1}$
eventually, so the ratios satisfy $x_{k+1}=1+1/x_k$ and converge to $\varphi$, giving $L=\varphi$, a
contradiction. Conversely $a_k=1$ eventually gives $q_{k+1}/q_k\to\varphi$.
\end{proof}

\subsection{The normalisation of \texorpdfstring{$\om$}{omega} and the threshold}\label{ss:norm}
For $\xi\in\Zt$ and $r\in\Q$ put $H(r)=\max(|\mathrm{num}|,\mathrm{den})$ as above, and let
$\om(\xi)$ be the supremum of the real $w$ for which
\begin{equation}\label{eq:omdef}
  |\xi-r|_2\le H(r)^{-(1+w)}
\end{equation}
holds for infinitely many $r\in\Q$. This is Koksma's $w_1$ in the $p$-adic normalisation of
\cite[Ch.~9]{Bugeaud04}: for $P(X)=\mathrm{den}\cdot X-\mathrm{num}$ one has
$|P(\xi)|_2=|\mathrm{den}|_2\,|\xi-r|_2$, which equals $|\xi-r|_2$ when the denominator is odd. A
pigeonhole over odd denominators $q\le H$ shows $\om(\xi)\ge1$ for every $\xi\in\Zt$, so
$\om=1$ is the generic value and the transcendence threshold is $\om>1$, not $\om>0$. The threshold
is the following, which we quote in its \emph{single-prime} form.\footnote{Do not quote the
\emph{several-primes} form of Ridout's theorem that circulates in some secondary sources, namely:
given $f\in\Z[X]$ of degree $\ge2$, let $\alpha$ be a \emph{real} root of $f$ and $\alpha_i$ be
$p_i$-adic roots ($p_1,\dots,p_t$ distinct primes); then for every $\varepsilon>0$ the inequality
$\min(1,|\alpha-A/B|_\infty)\prod_i\min(1,|\alpha_i-A/B|_{p_i})<\max(|A|,|B|)^{-2-\varepsilon}$ has
only finitely many coprime integer solutions. Its hypothesis that $f$ have a real root is a genuine
restriction on algebraic elements of $\Zt$, so the $t=1$ specialisation is not available in general.
Witness: $f(X)=X^{4}-3X^{3}+2X+6$ has $0$ real roots by Sturm's theorem in exact rational
arithmetic, no rational root, and no factorisation into integer quadratics (exhaustive over
$B\mid6$, $|A|\le40$, which by Gauss's lemma exhausts factorisations over $\Q$ of a monic quartic),
hence is irreducible over $\Q$ of degree $4$; while $f(1)=6\equiv0\pmod 2$ and $f'(1)=-3$ is odd, so
Hensel's lemma gives a root $\xi\in\Zt$, with $\xi\equiv16095319\pmod{2^{24}}$ and
$f(\xi)\equiv0\pmod{2^{800}}$, verified exactly. An exhaustive search over a small coefficient box
produced $1138$ such quartics, so this is generic rather than exotic.}

\begin{theoremR}[{\cite[Theorem 1.3]{BK18}, attributed there to Ridout \cite{Ridout58}}]
Let $\xi$ be a $p$-adic number and $\varepsilon>0$. Suppose there is a sequence
$(x_n/y_n)_{n\ge1}$ of rationals with $\gcd(x_n,y_n)=1$,
$2\le|x_1,y_1|<|x_2,y_2|<\cdots$, and
\[
  0<\Bigl|\xi-\frac{x_n}{y_n}\Bigr|_p<|x_n,y_n|^{-2-\varepsilon}\qquad(n=1,2,\dots),
\]
where $|x,y|:=\max(|x|,|y|)$ is the height of $x/y$. Then $\xi$ is transcendental.
\end{theoremR}

Theorem R is stated for an \emph{arbitrary} $p$-adic number, with our height normalisation and no
hypothesis on the minimal polynomial, and its hypothesis is exactly the family of approximants
produced in Step~4 of \S\ref{sec:main}. Equivalently $\om(\xi)\le1$ for every algebraic irrational
$\xi\in\Qt$. Appendix~\ref{app:schlickewei} re-derives it from Schlickewei's $p$-adic subspace
theorem; that derivation is not new and nothing depends on it.

Finally, the irrationality input, which Theorem R cannot supply: a rational $r$ has
$\om(r)=\infty$, so $\om>1$ does not by itself exclude rationality.

\begin{proposition}[{\cite[Theorem 10.1 and Corollary 10.5]{DJirr}}]\label{prop:irr}
$\PH(s)\notin\Q$ for every mechanical word $s$ of every irrational slope and every intercept.
\end{proposition}

Prior partial coverage: \cite[Theorem 2.7(b)]{MY04} gives this whenever the lower ones-density is
$<\beta$. We do not use \cite[Theorem 1]{LS21}.

\section{Proof of the main theorems}\label{sec:main}

\subsection{Proof of Theorem~\ref{thm:main}}

\emph{Step 1: choice of the exponent.} By Proposition~\ref{prop:icefloor},
$\ice(c_\gamma)\ge1+\varphi>2$. We must produce a fixed real $e>\max(2\Av,2)$ such that
arbitrarily long prefixes of $c_\gamma$ have prefix power at least $e$. Two cases, treated
separately because $\ice$ may be infinite.

\begin{itemize}
\item \emph{Case L ($\ice(c_\gamma)=\infty$).} By Proposition~\ref{prop:icefla} this happens exactly
  when the partial quotients of $\gamma$ are unbounded. Here \emph{every} real $e$ is admissible:
  since $\ice$ is a limit superior of prefix powers, for each $e$ there are arbitrarily long
  prefixes of $c_\gamma$ whose prefix power exceeds $e$. Fix, for definiteness,
  $e:=2\Av+1$. At the end of Step~5 we let $e\to\infty$; see \S\ref{sec:liouville}.
\item \emph{Case F ($\ice(c_\gamma)<\infty$).} By \eqref{eq:T} and $\ice\ge1+\varphi>2$ we have
  $\max(2\Av,2)<\ice(c_\gamma)$, so we may fix a real $e$ with $\max(2\Av,2)<e<\ice(c_\gamma)$.
\end{itemize}

No perturbation of $\ice$ is needed, and none is taken: the limit-superior definition supplies the
prefixes directly.

\emph{Step 2: the family of shadows, and their heights.} Replacing each such prefix by its
primitive root if necessary --- which only increases the prefix power, since $W=V^{t}$ gives $V$
prefix power $t$ times that of $W$ --- choose an infinite sequence $W_1,W_2,\dots$ of primitive
prefixes of $c_\gamma$ with $\ell_j:=|W_j|\to\infty$, $\ell_j\ge2$, and
\begin{equation}\label{eq:depth}
  \lcp\bigl(c_\gamma,W_j^{\infty}\bigr)\;\ge\;e\,\ell_j .
\end{equation}
The lengths of the primitive roots do tend to infinity: otherwise $c_\gamma$ would agree with one
fixed periodic word to unbounded depth, contradicting its aperiodicity. By
Definition~\ref{def:ice} the relation \eqref{eq:depth} involves no floor and no $O(1)$ loss.

Since $e>2$, Proposition~\ref{prop:bhz32} applies: each $W_j$ is a conjugate of a standard word, so
$W_j^{\infty}$ is balanced. Also $k_j:=|W_j|_1\ge1$ once $\ell_j$ is large, a Sturmian word of
irrational slope not being $0^{\infty}$; so Proposition~\ref{prop:shadow} gives
\begin{equation}\label{eq:Rj}
  R_j:=\PH\bigl(W_j^{\infty}\bigr)=\frac{c_{W_j}}{2^{\ell_j}-3^{k_j}}\in\Q,
  \qquad\text{denominator odd},
\end{equation}
and Proposition~\ref{prop:height} gives
\begin{equation}\label{eq:hbound}
  \log_2 H(R_j)\;<\;\Av\,\ell_j+\log_2(3\ell_j)+\log_2 3 .
\end{equation}
This bound holds \emph{at every level used}, not merely along the convergent family: its only input
is that $W_j^{\infty}$ is balanced, which Proposition~\ref{prop:bhz32} supplies. (We verified both
$c_W\le3\ell\max(2^{\ell},3^{k})$ and \eqref{eq:hbound} directly at every level used, for the
prefixes of $c_\gamma$ specifically; see Appendix~\ref{app:verif}.)

\emph{Step 3: the approximation exponent.} By Proposition~\ref{prop:iso} and \eqref{eq:depth},
$\vtwo(\PH(c_\gamma)-R_j)=\lcp(c_\gamma,W_j^{\infty})\ge e\ell_j$, so
$|\PH(c_\gamma)-R_j|_2\le2^{-e\ell_j}$. Combining with \eqref{eq:hbound},
$|\PH(c_\gamma)-R_j|_2\le H(R_j)^{-(1+w_j)}$ where
\begin{equation}\label{eq:wj}
  1+w_j=\frac{e\,\ell_j}{\log_2 H(R_j)}
  >\frac{e\,\ell_j}{\Av\,\ell_j+\log_2(3\ell_j)+\log_2 3}
  \;\xrightarrow[j\to\infty]{}\;\frac{e}{\Av}.
\end{equation}
Since $e>2\Av$, the limit exceeds $2$, so $w_j>1$ for all large $j$.

\emph{Step 4: the shadows are infinitely many, pairwise distinct, and distinct from the target.} If
$R_j=\PH(c_\gamma)$ then $\PH(W_j^{\infty})=\PH(c_\gamma)$, so by injectivity
(Proposition~\ref{prop:iso}) $c_\gamma=W_j^{\infty}$ would be periodic --- impossible, $\gamma$
being irrational. Hence $R_j\ne\PH(c_\gamma)$ for every $j$. Next, $R_j$ determines $W_j^{\infty}$
by injectivity, and $W_j$ primitive means $W_j^{\infty}$ has minimal period exactly $\ell_j$; since
$\ell_j\to\infty$, the $R_j$ take infinitely many distinct values. Passing to a subsequence they are
pairwise distinct with $H(R_j)\to\infty$, only finitely many rationals having bounded height. In
particular no shadow has $H=1$, where \eqref{eq:omdef} would be vacuous.

\emph{Step 5: the threshold, and excluding rationality.} By Steps~3 and~4 we have infinitely many
pairwise distinct rationals $R_j$, of strictly increasing height after passing to a subsequence,
with odd denominators, satisfying $|\PH(c_\gamma)-R_j|_2\le H(R_j)^{-(1+w_j)}$ and $w_j>1$. Hence
\begin{equation}\label{eq:omlb}
  \om\bigl(\PH(c_\gamma)\bigr)\;\ge\;\liminf_j w_j\;\ge\;\frac{e}{\Av}-1\;>\;1 .
\end{equation}
Choosing $\varepsilon>0$ with $1+w_j>2+\varepsilon$ for large $j$, Theorem R applies and
$\PH(c_\gamma)$ is \emph{not} an algebraic irrational.

It remains to exclude rationality, which Theorem R cannot do: a rational $r$ has $\om(r)=\infty$.
By Proposition~\ref{prop:irr}, $\PH(c_\gamma)\notin\Q$ --- either directly, $c_\gamma$ being a
mechanical word of slope $\gamma$ and intercept $\{\gamma\}$, or from $\PH(1c_\gamma)\notin\Q$
through Step~6: if $\PH(c_\gamma)$ were rational then so would be
$\PH(1c_\gamma)=(2\PH(c_\gamma)-1)/3$. There is no circularity, Step~6 transferring rationality by
an invertible $\Q$-affine map independently of the transcendence conclusion. Therefore
$\PH(c_\gamma)$ is transcendental.

\emph{Step 6: transfer along the shift orbit.} Specialising \eqref{eq:affinegen} to $m=1$ with
$v=1c_\gamma$ (so $k_1=1$, $c_1=1$) and with $v=0c_\gamma$ (so $k_1=0$, $c_1=0$), and using
$\sigma(1c_\gamma)=\sigma(0c_\gamma)=c_\gamma$:
\begin{equation}\label{eq:affine}
  2\,\PH(c_\gamma)=3\,\PH(1c_\gamma)+1=\PH(0c_\gamma),
  \qquad\text{equivalently}\qquad \PH(1c_\gamma)=\frac{2\PH(c_\gamma)-1}{3}.
\end{equation}
All three maps are $\Q$-affine and invertible over $\Q$, so the three values are simultaneously
algebraic or simultaneously transcendental; the same holds for every $\PH(\sigma^{m}v)$ by
\eqref{eq:affinegen}. More generally \eqref{eq:affinegen} shows that $\PH$ of any two
\emph{tail-equivalent} words are $\Q$-affinely related, through
$3^{k_m(v)}\PH(v)-3^{k_m(w)}\PH(w)=c_m(w)-c_m(v)$. This proves Theorem~\ref{thm:main}. \qed

\medskip
The identity \eqref{eq:affine} is not new. \cite[Lemma 37]{LS21} (equivalently
\cite[Lemma 25]{LS21}) gives $\PH_{\R}(0c_\alpha)=3\PH_{\R}(1c_\alpha)+1$ in the \emph{real}
normalisation for $\alpha>\beta$, and the $2$-adic form used here, valid for every irrational
slope, is the author's prior work \cite[Corollaries 7.1 and 8.8]{DJirr}. What this paper adds is
only the observation that, being $\Q$-affine, it carries transcendence along the shift orbit. We verified
\eqref{eq:affine} exactly modulo $2^{400}$ at three slopes (Appendix~\ref{app:verif}).

\subsection{Proof of Theorem~\ref{thm:quant}}
Steps 1--5 give $\om(\PH(c_\gamma))\ge e/\Av-1$ for every admissible $e$. In Case F, $e$ may be any
real number below $\ice(c_\gamma)$; in Case L, any real number at all. Taking the supremum over
admissible $e$ gives $\om(\PH(c_\gamma))\ge\ice(c_\gamma)/\Av-1$, which together with
Proposition~\ref{prop:icefla} is the assertion. \qed

\subsection{The margin is a fixed \texorpdfstring{$\delta$}{delta}}
\begin{lemma}\label{lem:margin}
Assume \eqref{eq:T} and set $\kappa:=\ice(c_\gamma)-2\Av>0$. In Case F the value
$e:=2\Av+\kappa/2$ is admissible in Step~1, and then $w_j\to1+\delta$ with
\[
  \delta=\frac{\kappa}{2\Av}>0,
\]
independent of $j$, of $\ell_j$ and of $H(R_j)$. Moreover the loss in \eqref{eq:wj} is
$O(\log\ell_j/\ell_j)=O(\log\log H_j/\log H_j)\to0$, so there is an explicit
$L_0=L_0(\kappa,\gamma)$ with $w_j\ge1+\delta/2>1$ whenever $\ell_j\ge L_0$. In Case L,
$\kappa=\infty$ and every finite $\delta$ is achieved.
\end{lemma}

\begin{proof}
$e=2\Av+\kappa/2$ lies strictly between $\max(2\Av,2)$ and $\ice=2\Av+\kappa$, so it is admissible;
the limit in \eqref{eq:wj} is $e/\Av=2+\kappa/(2\Av)$. For the loss, \eqref{eq:hbound} gives
$1+w_j\ge(e/\Av)\bigl(1-O(\log\ell_j/\ell_j)\bigr)$, and $\log_2 H_j\asymp\Av\ell_j$.
\end{proof}

This is the quantitative point on which the argument turns. For an algebraic irrational the excess
of $\om$ over $1$ \emph{is} a vanishing slack: for a quadratic irrational, Liouville's inequality
caps it at $1+\log_2 9/\log_2 H\to1$ (Appendix~\ref{app:verif}). Here the excess is the constant
$\delta$ while the slack is $O(\log\log H/\log H)$. A vanishing slack cannot eat a fixed $\delta$.

\subsection{Proof of Corollary~\ref{cor:uncond}}
By Proposition~\ref{prop:icefloor}, $\ice(c_\gamma)\ge1+\varphi$. If $\gamma<\gs=(1+\varphi)/(2\log_2 3)$
then $\gamma\log_2 3<(1+\varphi)/2$; also $1<(1+\varphi)/2$ since $\varphi>1$. Hence
$\Av=\max(1,\gamma\log_2 3)<(1+\varphi)/2$, i.e.\ $2\Av<1+\varphi\le\ice(c_\gamma)$, so \eqref{eq:T}
holds with strict inequality and Theorem~\ref{thm:main} applies. At $\gamma=\beta=\log_3 2$ one has
$\gamma\log_2 3=1$, so $\Av=1$ and $2\Av=2<2.618\le\ice$. \qed

\begin{remark}[Threshold arithmetic]\label{rem:threshold}
Explicitly, $(1+\varphi)/\Av>2\iff\Av<(1+\varphi)/2=(3+\sqrt5)/4=1.3090169943749474241\ldots$. For
$\gamma\le\beta$ one has $\Av=1$ and the condition is automatic; for $\gamma>\beta$ it is
$\gamma\log_2 3<(3+\sqrt5)/4$, i.e.\ $\gamma<\gs$. The constant $\Av$ is the right height exponent
across the whole range: for $\gamma\le\beta$ the denominator $|2^{\ell}-3^{k}|\asymp2^{\ell}$
dominates and $\log_2 H\asymp\ell$; for $\gamma>\beta$, $3^{k}\asymp2^{\gamma\ell\log_2 3}$ dominates
and $\log_2 H\asymp\gamma\ell\log_2 3$. Both regimes are confirmed numerically
(Appendix~\ref{app:verif}).
\end{remark}

\begin{remark}[The two $\Xi$ conventions; read once, not used again]\label{rem:xi}
Two conventions for a symbol $\Xi$ are in circulation in this circle of papers, and they denote
\emph{negatives of each other}. In \cite[Proposition 3.2]{DJirr} the definition carries no sign,
$\Xi_\alpha:=\sum_{j\ge0}3^{-(j+1)}2^{\lfloor j\alpha\rfloor}$, and then
$\PH(1c_\beta)=-\Xi_\alpha$ with $\alpha=\log_2 3$, $\beta=\log_3 2$. In
\cite[Theorem 4.1]{DJedge} the sign is inside the definition,
$\Xi_{\alpha,b}:=-\frac13\sum_{i\ge0}3^{-i}2^{S_i}=-\sum_{i\ge0}3^{-i-1}2^{\lfloor i\alpha+b\rfloor}$
with $S_i=\lfloor i\alpha+b\rfloor$, and then $\PH(1c_\beta)=+\Xi_{\alpha,0}$. Equivalently
$\Xi_{\alpha,0}=-\Xi_\alpha$. Both papers are internally correct; crossing the conventions is not.
We verified exactly, with $\PH$ computed from the parity-vector definition alone and certified
floors, that $\PH(1c_\beta)=-\Xi_\alpha$ and $\PH(1c_\beta)=+\Xi_{\alpha,0}$ both hold modulo
$2^{3000}$ (1893 terms summed) while the opposite signs both fail, and that
$\vtwo(\PH(1c_\beta))=0$. To avoid the collision, no statement in this paper uses $\Xi$; everything
is phrased in terms of $\PH(c_\gamma)$, $\PH(1c_\gamma)$ and $\PH(0c_\gamma)$.
\end{remark}

\section{The parity obstruction: why \texorpdfstring{$c_\gamma$}{c-gamma} and not
\texorpdfstring{$1c_\gamma$}{1c-gamma}}\label{sec:parity}

It is natural to try to run \S\ref{sec:main} directly on $1c_\gamma$, the word whose $\PH$-value is
the object of \cite{DJirr} and of the realizer problem of \cite{DJedge}. That fails, and the failure
is not an artifact. By the depth law of \cite[Theorems 5.3 and 8.3]{DJirr} --- the author's prior work, and the
one place where both convergent parities are computed exactly --- the length-$q_n$ prefix of
$1c_\gamma$ has prefix power exactly $1$ at every \emph{even} $n$, whereas at odd $n$ it is
$(q_n+q_{n+1}-1)/q_n$. Consequently
\begin{equation}\label{eq:parity}
  \ice(1c_\gamma)=1+\limsup_{n\ \mathrm{odd}}\frac{q_{n+1}}{q_n},
  \qquad\text{against}\qquad
  \ice(c_\gamma)=1+\limsup_{k}\frac{q_{k+1}}{q_k},
\end{equation}
and the first can be strictly smaller. Both halves of \eqref{eq:parity} were verified
independently at every level of four slopes (Appendix~\ref{app:verif}).

That the gap can be fatal is shown by an explicit witness.

\begin{example}[The $1c_\gamma$ route genuinely fails]\label{ex:witness}
Take
\[
  \gamma=[0;2,1,4,1,8,1,16,1,32,1,\dots]=0.35329542\ldots<\gs,
\]
so that $a_n\to\infty$ along odd $n$ while $a_n=1$ for even $n$. For odd $n$ one has $a_{n+1}=1$
and $q_{n-1}/q_n\to0$, so the odd-$n$ ratios collapse to $1$: we computed
\[
  q_{n+1}/q_n=1.50000,\ 1.21429,\ 1.11333,\ 1.05918,\ 1.03035,\ 1.01539,\ 1.00775,\ 1.00389,\ 1.00195
\]
at $n=1,3,5,\dots,17$. Hence $\ice(1c_\gamma)=2$ \emph{exactly}, the $1c_\gamma$ route yields only
$\om\ge1$ --- the threshold of Theorem R, with zero margin, proving nothing --- and the measured
certified lower bounds for $\om$ along the levels decrease, $1.1429$, $1.1067$, $1.0588$, towards
$1$. On the \emph{same} slope the full limit superior over all $k$ is infinite (ratios
$1.50$, $4.67$, $1.21$, $8.82$, $1.11$, $16.90$, $1.06$, $32.94$, $1.03$), so
$\ice(c_\gamma)=\infty$; this is Case~L of Step~1, and the $c_\gamma$ route gives a measured
$\om\ge16.886$ at the levels computed.
\end{example}

So Theorem~\ref{thm:main} is formulated for $c_\gamma$ of necessity, and Step~6 is what carries the
conclusion back to $1c_\gamma$ and $0c_\gamma$ through \eqref{eq:affine}. The transfer is lossless
in the limit: a shadow $R=u/v$ of $\PH(c_\gamma)$ with depth $D$ and odd $v$ becomes $2R=2u/v$ for
$0c_\gamma$ and $(2R-1)/3=(2u-v)/(3v)$ for $1c_\gamma$, each with depth exactly $D+1$, each with
odd denominator, and with height cost only $\log_2 H+1$, respectively $\log_2 H+\log_2 3$. Since
$\log_2 H(R_j)\asymp\Av\ell_j\to\infty$, the limit in \eqref{eq:wj} is unchanged for all three
words. In our sweep over every prefix length $\ell<900$ of $1c_\gamma$, no other prefix family repairs
the parity loss.

\section{Why the truncations fail and the periodic approximants succeed}\label{sec:truncations}

The question of approximating $\PH(1c_\beta)$ at the resonance slope was first posed in
\cite{DJedge}, where the natural approximants are the Hensel truncations of the defining series
rather than the periodic shadows used here. It is worth recording exactly why that family cannot
prove a statement of this kind, since the contrast isolates what the shadows supply.

Fix $\alpha=\log_2 3$, $\beta=\log_3 2$ and $S_i=\lfloor i\alpha\rfloor$, and write
$\PH(1c_\beta)=-\frac13\sum_{i\ge0}3^{-i}2^{S_i}$ in the convention of
\cite[Theorem 4.1]{DJedge} (Remark~\ref{rem:xi}). Its $N$-term truncation is
\[
  \Xi_N=-\frac{C_N}{3^{N}},\qquad C_N:=\sum_{i<N}3^{\,N-1-i}2^{S_i}\in\Z_{>0},
\]
so $P_N:=-C_N$ and $Q_N:=3^{N}$. The fraction is already reduced: the $i=N-1$ term is $2^{S_{N-1}}$
and every other term is divisible by $3$, so $3\nmid C_N$.

\begin{proposition}[Determinant identity]\label{prop:det}
For every $N\ge1$,\; $P_NQ_{N+1}-P_{N+1}Q_N=+3^{N}2^{S_N}$.
\end{proposition}

\begin{proof}
From the definition, $C_{N+1}=3C_N+2^{S_N}$. Hence
$P_NQ_{N+1}-P_{N+1}Q_N=(-C_N)3^{N+1}+C_{N+1}3^{N}=3^{N}(C_{N+1}-3C_N)=3^{N}2^{S_N}$.
\end{proof}

\begin{remark}[The determinant is an $S$-unit, and the product formula prices it]\label{rem:sunit}
Write $D_N=3^{N}2^{S_N}$. Then $|D_N|_2|D_N|_3|D_N|_\infty=2^{-S_N}\cdot3^{-N}\cdot3^{N}2^{S_N}=1$
exactly: $D_N$ is an $S$-unit for $S=\{2,3,\infty\}$, so its $2$-adic depth is paid, to the last
bit, by its archimedean size. Consequently
$\vtwo(D_N)/\log_2|D_N|=S_N/(N\log_2 3+S_N)\to\frac12$, the exponent a Minkowski-generic lattice
vector already achieves. No linear combination of two consecutive truncations can therefore cross
the threshold; the obstruction is a conservation law, not a loose estimate.
\end{remark}

\begin{proposition}[Truncations are subcritical against the Weil height]\label{prop:trunc}
$S_N-\log_2 Q_N\in(-1,0]$ for every $N$, while $S_N-\log_2 H(\Xi_N)=-\log_2 N+O(1)$.
\end{proposition}

\begin{proof}
Since $2^{\alpha}=3$ we have $2^{S_i}=2^{i\alpha-\{i\alpha\}}=3^{i}2^{-\{i\alpha\}}$, so each term of
$C_N$ equals $3^{N-1}2^{-\{i\alpha\}}\in(\tfrac12 3^{N-1},3^{N-1}]$ --- this is \eqref{eq:terms}
again. Summing the $N$ terms, $\tfrac12N3^{N-1}<C_N\le N3^{N-1}$, so
$H(\Xi_N)=\max(C_N,3^{N})=C_N$ once $N\ge4$ and
$\log_2 H(\Xi_N)=N\log_2 3+\log_2 N-\log_2 3+O(1)$. Meanwhile $S_N=N\alpha-\{N\alpha\}$ gives
$S_N-\log_2 Q_N=-\{N\alpha\}\in(-1,0]$, and subtracting the two displays gives the second claim.
\end{proof}

Both propositions are \textbf{T1}; both were checked in exact integer arithmetic
(Appendix~\ref{app:verif}). The contrast with \S\ref{sec:main} is now explicit. A truncation has
depth $S_N\approx N\log_2 3$ against height $\approx N\log_2 3+\log_2 N$, so its exponent tends to
$1$ --- it is \emph{sub}critical, and it is subcritical because the Weil height, not the
denominator, is what Theorem R and the $p$-adic Liouville inequality use. A shadow $\PH(W^{\infty})$
has depth at least $e\,|W|$ against height at most $\Av|W|+O(\log|W|)$, so its exponent tends to
$e/\Av>2$. The gain is in the depth and not in the height: a Sturmian prefix repeats, with initial
power at least $1+\varphi$ by Proposition~\ref{prop:icefloor}, whereas a truncation of a series
repeats not at all. That is the whole of the difference.

\section{Unbounded partial quotients: the Liouville alternative}\label{sec:liouville}

\begin{theorem}\label{thm:liouville}
Let $\gamma\in(0,1)$ be irrational with unbounded partial quotients. Then
$\om(\PH(c_\gamma))=\infty$; that is, $\PH(c_\gamma)$ is a $2$-adic Liouville number, a $p$-adic
$U_1$-number in Mahler's classification. The same holds for $\PH(1c_\gamma)$, $\PH(0c_\gamma)$ and
every shift. In particular no finite irrationality measure exists for these numbers.
\end{theorem}

\begin{proof}
Unbounded partial quotients give $\limsup_k q_{k+1}/q_k=\infty$, hence $\ice(c_\gamma)=\infty$ by
Proposition~\ref{prop:icefla}; Theorem~\ref{thm:quant} then gives $\om(\PH(c_\gamma))=\infty$. This
is Case~L of Step~1: every real $e$ is admissible, and \eqref{eq:omlb} holds for each, so
$\om\ge e/\Av-1$ for every $e$. For the other words, $\om$ is unchanged up to the $O(1/\log H)$
cost of the transfer computed at the end of \S\ref{sec:parity}, which does not affect its being
infinite. A $2$-adic number with $\om=\infty$ attained along infinitely many distinct rationals is
a $U_1$-number by the $p$-adic Mahler classification \cite[Ch.~9]{Bugeaud04}.
\end{proof}

Theorem~\ref{thm:liouville} refines the corresponding statement of the author's earlier notes, where
the limit superior was taken over odd $n$ only; passing from $1c_\gamma$ to $c_\gamma$ removes the
parity restriction, so unbounded partial quotients at \emph{either} parity now suffice.

For $\gamma=\beta=\log_3 2$ it is unknown whether the partial quotients of $\log_2 3$ are bounded.
We computed $1456$ certified partial quotients, with maximum $3308$; the largest certified lower
bound on $\om(\PH(c_\beta))$ obtainable from a single level in our range is $23.274$. Whether
$\om(\PH(c_\beta))$ is finite is therefore open, and equivalent to the boundedness of the partial
quotients of $\log_2 3$. \emph{Transcendence does not depend on it}: that is the content of
Corollary~\ref{cor:uncond}.

\section{Bounded partial quotients: a finite effective measure}\label{sec:measure}

This section is logically independent of \S\ref{sec:main} and is stated at a narrower scope than the
rest of the paper. \textbf{It was not covered by the independent verification reported in
Appendix~\ref{app:verif}}, whose remit was Theorems~\ref{thm:quant}--\ref{thm:main}; the numerics
quoted below are the author's own.

We first record the abstract mechanism.

\begin{lemma}[Lower bound for shadow heights]\label{lem:hlow}
Let $R\ne R'$ be rationals with odd denominators. Then
$\vtwo(R-R')\le\log_2 H(R)+\log_2 H(R')+1$.
\end{lemma}

\begin{proof}
Write $R=A/B$, $R'=A'/B'$ in lowest terms with $B,B'$ odd. Then $R-R'=(AB'-A'B)/(BB')$ with
$BB'$ odd, so $\vtwo(R-R')=\vtwo(AB'-A'B)$, and $AB'-A'B$ is a nonzero integer of absolute value at
most $|A|B'+|A'|B\le2H(R)H(R')$.
\end{proof}

\begin{lemma}[Abstract measure lemma]\label{lem:P}
Let $\xi\in\Zt$ be irrational and let $(R_k)_{k\ge1}$ be pairwise distinct rationals with odd
denominators, $L_k:=\vtwo(\xi-R_k)<\infty$ and $h_k:=\log_2 H(R_k)\to\infty$ strictly increasing.
Put
\[
  \theta:=\liminf_k\frac{L_k}{h_k},\qquad \Theta:=\limsup_k\frac{L_k}{h_k},\qquad
  g:=\limsup_k\frac{h_{k+1}}{h_k}.
\]
Assume $\Theta<\infty$, $g<\infty$ and $\theta>1$. Then for every
$\mu>\max\bigl(\Theta,\,1+g/(\theta-1)\bigr)$ there is $C>0$ with
$|\xi-r|_2\ge C\,H(r)^{-\mu}$ for all $r\in\Q\cap\Zt$.
\end{lemma}

\begin{proof}
\emph{Basic step.} Let $r=p/q$ in lowest terms with $q>0$ odd, $H=H(r)$, and suppose $r\ne R_k$,
$R_k=A_k/B_k$. By Lemma~\ref{lem:hlow} applied to $r$ and $R_k$,
$\vtwo(R_k-r)\le h_k+\log_2 H+1$. If moreover $L_k>h_k+\log_2 H+1$ then
$\vtwo(\xi-R_k)=L_k>\vtwo(R_k-r)$, so the two valuations differ and the ultrametric equality gives
$\vtwo(\xi-r)=\vtwo(R_k-r)\le h_k+\log_2 H+1$, whence
\begin{equation}\label{eq:basic}
  |\xi-r|_2\ \ge\ 2^{-1}\,2^{-h_k}H^{-1}.
\end{equation}

\emph{Measure.} Fix $\mu$ as stated and choose $\epsilon>0$ with
$\mu>\max\bigl(\Theta+\epsilon,\,1+(g+\epsilon)/(\theta-1-\epsilon)\bigr)$, then $N$ with
$\theta-\epsilon<L_k/h_k<\Theta+\epsilon$ and $h_{k+1}/h_k<g+\epsilon$ for all $k\ge N$.

If $r=R_k$ for some $k\ge N$ then $H=2^{h_k}>1$ and
$|\xi-r|_2=2^{-L_k}=H^{-L_k/h_k}\ge H^{-(\Theta+\epsilon)}\ge H^{-\mu}$. The finitely many $k<N$
have bounded height and are absorbed into $C$.

Otherwise $r\ne R_k$ for all $k$. Put $X=\log_2 H$ and let $k$ be the largest index with
$h_k\le(\mu-1)X-1$; for $X$ large this $k$ exists and $k\ge N$. Then $h_{k+1}>(\mu-1)X-1$, so
$h_k>h_{k+1}/(g+\epsilon)>\bigl((\mu-1)X-1\bigr)/(g+\epsilon)$. Now
$L_k>(\theta-\epsilon)h_k$, and $L_k>h_k+X+1$ holds as soon as
$h_k(\theta-1-\epsilon)>X+1$, for which it suffices that
$(\theta-1-\epsilon)\bigl((\mu-1)X-1\bigr)>(g+\epsilon)(X+1)$; as $X\to\infty$ this is implied by
$(\theta-1-\epsilon)(\mu-1)>g+\epsilon$, i.e.\ by the choice of $\epsilon$. So the basic step
applies, and \eqref{eq:basic} together with $h_k\le(\mu-1)X-1$ gives
$|\xi-r|_2\ge2^{-1-h_k-X}\ge2^{-\mu X}=H^{-\mu}$. The finitely many $r$ with $X$ below the
resulting threshold have bounded height and are absorbed into $C$.
\end{proof}

The hypothesis $\Theta<\infty$ is not a convenience: $\Theta=\infty$ says exactly that $\xi$ is a
$2$-adic Liouville number, i.e.\ that the conclusion is false. Lemma~\ref{lem:P} therefore fails
precisely where it must.

\begin{theorem}\label{thm:Bstar}
Let $\gamma\in(0,\beta]$ be irrational with bounded partial quotients, $\beta=\log_3 2$. Put
\[
  \lambda:=\liminf_k\frac{q_{k+1}}{q_k},\qquad \Lambda:=\limsup_k\frac{q_{k+1}}{q_k}
  \quad(\text{both finite, and }\lambda>1).
\]
Then for every $\mu>\max\bigl(1+\Lambda,\,1+\Lambda/\lambda\bigr)$ there is an effective $C>0$ with
\[
  \bigl|\PH(c_\gamma)-r\bigr|_2\ \ge\ C\,H(r)^{-\mu}\qquad\text{for all }r\in\Q\cap\Zt,
\]
and the same with $\PH(1c_\gamma)$ or $\PH(0c_\gamma)$ in place of $\PH(c_\gamma)$, after adjusting
$C$. For a constant partial quotient $a$, with $\lambda=\Lambda=(a+\sqrt{a^{2}+4})/2$, this reads
$\mu>1+\lambda$; e.g.\ $\mu>\varphi^{2}=2.6180\ldots$ for $\gamma=1/\varphi$ and
$\mu>2+\sqrt2=3.4142\ldots$ for $\gamma=\sqrt2-1$.
\end{theorem}

\begin{proof}
Let $W_k$ be the prefix of $c_\gamma$ of length $q_k$ and $R_k=\PH(W_k^{\infty})$; these are
pairwise distinct with odd denominators, as in Step~4. By Proposition~\ref{prop:icefla} and the
proof of \cite[\S4.2]{BHZ06}, $L_k:=\vtwo(\PH(c_\gamma)-R_k)=q_{k+1}+q_k-2$. Since
$\gamma\le\beta$ we have $\Av=1$, so \eqref{eq:hbound} gives the upper bound
$h_k<q_k+\log_2(3q_k)+\log_2 3$. For the lower bound, Lemma~\ref{lem:hlow} applied to $R_k\ne
R_{k+1}$ gives $L_k=\vtwo(R_k-R_{k+1})\le h_k+h_{k+1}+1$ (the valuation of the difference being
$\min(L_k,L_{k+1})=L_k$), whence
\[
  h_k\ \ge\ L_k-1-h_{k+1}\ >\ (q_{k+1}+q_k-3)-q_{k+1}-\log_2(3q_{k+1})-\log_2 3
  \ =\ q_k-O(\log q_{k+1}).
\]
Bounded partial quotients give $q_{k+1}\le(M+1)q_k$, so $\log q_{k+1}=O(\log q_k)$ and therefore
$h_k=q_k+O(\log q_k)$. Consequently
\[
  \theta=1+\lambda,\qquad \Theta=1+\Lambda,\qquad g=\Lambda,
\]
all finite, and $\theta>1$; indeed $q_{k+1}/q_k=a_{k+1}+q_{k-1}/q_k\ge1+1/(M+1)$, so
$\lambda>1$. Lemma~\ref{lem:P} applies with
$\max(\Theta,1+g/(\theta-1))=\max(1+\Lambda,1+\Lambda/\lambda)$. For a constant partial quotient
$\lambda=\Lambda>1$ and the maximum is $1+\lambda$. The statement for $\PH(1c_\gamma)$,
$\PH(0c_\gamma)$ follows from \eqref{eq:affine}, which changes $H(r)$ by a bounded factor.
\end{proof}

\begin{remark}[Scope of Theorem~\ref{thm:Bstar}, honestly]\label{rem:Bscope}
The restriction $\gamma\le\beta$ is needed, and is narrower than one might expect. The upper bound
\eqref{eq:hbound} on $h_k$ is available at every slope, but Lemma~\ref{lem:P} also needs a
\emph{lower} bound, to make $\Theta$ and $g$ finite, and the only elementary source we have for one
is Lemma~\ref{lem:hlow}, which yields $h_k\ge q_k+q_{k+1}-3-h_{k+1}$. For $\Av=1$ this is
$q_k-O(\log q_{k+1})$ and suffices; for $\Av>1$ the right-hand side is negative and the argument
breaks. We do not know whether $\gcd(c_W,2^{\ell}-3^{k})$ can be large enough to make $h_k$ much
smaller than $\Av\,q_k$, and we make no claim for $\gamma\in(\beta,\gs)$.
\end{remark}

\section{Scope and open problems}\label{sec:scope}

Four limitations of the argument, with the open problems they carry: general intercepts
(\S\ref{ss:intercepts}, Question~\ref{q:intercepts}); extending the range towards all slopes
(\S\ref{ss:aboveg}, Question~\ref{q:subspace}); the integer target (\S\ref{ss:integer}); and the
real place (\S\ref{ss:realplace}).

\subsection{General intercepts}\label{ss:intercepts}
Theorem~\ref{thm:main} covers $c_\gamma$, $1c_\gamma$, $0c_\gamma$ and the whole shift orbit of
$c_\gamma$ --- equivalently its tail-equivalence class, the countably many intercepts
$\{k\gamma\}$. It does \emph{not} cover arbitrary intercepts, and this is an obstruction rather than
an omission. Every Sturmian word begins in infinitely many squares \cite[Theorem 1]{ADQZ2001}, so
$\ice\ge2$ always; but by \cite[Theorem 1.1]{BHZ06} there are slopes $\alpha$ carrying a Sturmian
word $\omega\in X_\alpha$ with $\ice(\omega)=2$ \emph{exactly} --- namely those for which, for each
pair $(s,t)$ with $s>1$, only finitely many $k$ have $(a_k,a_{k+1})=(s,t)$ or
$(a_k,a_{k+1},a_{k+2})=(1,1,t)$. For such a word the argument yields $\om\ge2/\Av-1=1$ at $\Av=1$:
the threshold of Theorem R with $\delta=0$, which proves nothing. The full generality of
\cite[Corollary 10.5]{DJirr} therefore does not transfer to transcendence, and no claim is made for
it. For the golden slope no such word exists, since $(a_k,a_{k+1},a_{k+2})=(1,1,1)$ recurs; for
$\gamma=\beta$ whether such a word exists is open, being a statement about recurrence of patterns in
the partial quotients of $\log_2 3$.

\begin{question}\label{q:intercepts}
Is $\PH(s)$ transcendental for every Sturmian word $s$ of every intercept, at least for
$\gamma<\gs$?
\end{question}

\subsection{Slopes above \texorpdfstring{$\gs$}{gamma*}}\label{ss:aboveg}
For $\gamma\ge\gs$ the condition \eqref{eq:T} becomes a genuine hypothesis on the partial quotients,
and Theorem~\ref{thm:main} is conditional. The boundary is sharp in our method: on noble slopes,
where $\ice=1+\varphi$ exactly, the limit of $\om$ along the levels is $(1+\varphi)/\Av-1$, which is
$>1$ for $\gamma=0.82200179\ldots$ and $<1$ for $\gamma=0.84889772\ldots$; both sides were verified
(Appendix~\ref{app:verif}). Improving the range requires either a better height exponent than $\Av$
or an input stronger than Theorem R at the single place $2$.

\begin{question}\label{q:subspace}
Can the range be extended towards all slopes by an $S$-unit or Subspace argument over
$\{\infty,2,3\}$, in place of the single place $2$?
\end{question}

\noindent
This is the route sketched by Cassidy \cite[\S4.6]{Cassidy26}: build integer points
from the shadow data and apply Schlickewei's $p$-adic Subspace Theorem --- Theorem S of
Appendix~\ref{app:schlickewei} --- at the three places $\{\infty,2,3\}$, rather than Theorem R at
the one place $2$. The slope-range bookkeeping is not carried out there, as that note itself
indicates, and it is what such a proof would have to supply. The cost of the extra coordinate is
explicit: with three forms the place-product must be pushed below $H(\mathbf{x})^{-3-\varepsilon}$,
where the single-prime route needs only $H(\mathbf{x})^{-2-\varepsilon}$ with two. Against that cost one may set a concrete gain. Working at the places $2$ and $3$ together puts
both primes in $S$, so the shadow denominators $2^{\ell}-3^{k}$ may be handled with both primes
available at once, rather than at the single place $2$. That is the obstruction identified in
\S\ref{ss:notcor} --- these denominators are not $S$-units for any fixed finite $S$, so the single
prime $2$ buys no unit of exponent --- and it is the one an argument over $\{\infty,2,3\}$ would
be in a position to attack.

\subsection{The integer target}\label{ss:integer}
A different and more demanding question is attached to the realizer problem of \cite{DJedge}: an
effective bound $|\cdot|_2\ge C\,y^{-\mu}$ against positive \emph{integers} $y$, which by
\cite[Proposition 7.2]{DJedge} is equivalent to a sublinear zero-run law on the Hensel digits.
Theorem~\ref{thm:liouville} does not obstruct it: the shadows all have large odd denominators ---
$2^{\ell}-3^{k}=\pm1$ only for $(\ell,k)\in\{(1,1),(2,1),(3,2)\}$, by an elementary argument ---
so none is an integer, and the Liouville blow-up does not descend to $\Z$. This remains open.

\subsection{The real place, and a remark of L\'opez and Stoll}\label{ss:realplace}
\begin{remark}\label{rem:real}
Equation \eqref{eq:terms} shows that at the resonance slope the terms of the defining series lie in
$(\frac16,\frac13]$ and do not tend to $0$: the real series \emph{diverges}, and the $2$-adic place
is the only one available. L\'opez and Stoll reach the same point from the other side: their real
staircase diverges at $x=\beta$. Their function $F(x):=\PH_{\R}(m_x)$ on $(\beta,1]$
\cite[Definition 26]{LS09} is a devil's staircase, and in \cite[\S4]{LS09} they remark: ``The
function is similar to other devil's staircases. \dots It seems that $F$ additionally maps
irrationals to transcendental numbers. We have no proof.'' That remark concerns the \emph{real}
function $F$, for slopes $>\beta$, and is explicitly unproven. The present paper proves
transcendence for the \emph{$2$-adic} $\PH$, on $(0,\gs)$, by a different mechanism. The two
statements are about different numbers and neither implies the other.
\end{remark}

\section{Related work}\label{sec:related}

L\'opez and Stoll are the origin of the explicit study of $\PH$ on Sturmian words. Their
\cite[Theorem 1]{LS09} is the closed $2$-adic series \eqref{eq:LS09series} and
\cite[Corollary 2]{LS09} the associated generalised continued fraction; \cite[Lemma 11]{LS09} is the
one-position formula for mechanical words, with \cite[Ch.~2]{Lothaire02} for the framework;
\cite[Lemma 12]{LS09} is the rational-slope shadow; \cite[\S4]{LS09} contains the irrationality of
the real value and the remark quoted in Remark~\ref{rem:real}; and \cite[Lemma 43, \S11]{LS21} is
the term-size estimate \eqref{eq:terms} with the distribution of the terms, while
\cite[Lemma 37]{LS21} is the real form of \eqref{eq:affine}. The approximation depths used
throughout this paper are the term exponents of their 2009 series. We do not rely on
\cite[Theorem 1]{LS21}.

Two contemporaneous independent contributions should be recorded, both concerning
\emph{irrationality}. We reuse the characterisations of \cite[\S1]{DJirr}. Pham
\cite[Theorem 3, Corollaries 4--5]{Pham26}, in an unrefereed note dated 22 July 2026, transported
Dubickas's counting argument to rational $2$-adic integers and drew the conclusion that every
Sturmian parity word has irrational $\PH$-value; his Theorem 3 is moreover sensitive to the drift
of the word. Cassidy \cite[Theorems R and S]{Cassidy26}, in an unrefereed note of September
2026, proved by a $2$-adic repetition (Liouville) criterion that no rational with odd denominator
has an eventually Sturmian parity vector, for every slope and intercept, and for a wider class of
maps: $3x+r$, $5x+1$ at slopes below $0.804$, and Mahler's map; his second route, through initial
squares, is the method of \cite[\S10]{DJirr}. Both were found independently of \cite{DJirr}, and
Cassidy's note was posted first. Neither note is refereed and we make no comparative claim about
either.

Transcendence of $\PH(s)$ for every Sturmian word $s$ was first proposed by Cassidy
\cite[\S4.6]{Cassidy26}, with a sketch via Schlickewei's $p$-adic Subspace Theorem, recorded
there as unaudited. The complete proof given here covers characteristic words of every irrational
slope $\gamma<\gs$ and goes by a different route, Ridout's theorem in the single-prime form
\cite[Theorem 1.3]{BK18} with the Berth\'e--Holton--Zamboni bound on initial powers; the general
claim remains open. See \S\ref{ss:aboveg} for what an argument over $\{\infty,2,3\}$ would
have to supply.


On the archimedean side the analogous programme is well developed: Ferenczi--Mauduit \cite{FM97},
Adamczewski--Bugeaud \cite{AB07}, and for Hecke--Mahler values Bugeaud--Laurent
\cite{BugeaudLaurent23}. As explained in \S\ref{ss:notcor}, the hypotheses of these results are on a
digit expansion and our object does not satisfy them.

\section*{Acknowledgments}
The author thanks Josefina L\'opez for identifying the earlier sources of several ingredients used
here --- the isometry, the one-position formula, the term-size bound of \cite[Lemma 43]{LS21}, and
the periodic approximants of \cite[Theorem 1]{LS09} --- and for the correction recorded in
Remark~\ref{rem:xi}.

Theorems~\ref{thm:quant}--\ref{thm:main} and their supporting computations were re-checked in a
separate AI session (Anthropic's Claude), working blind from the primary sources before reading the
proofs, with its own exact-arithmetic scripts. It re-established the three primary inputs from those
sources rather than from this paper: Ridout's exponent $2+\varepsilon$ with the height
$\max(|\mathrm{num}|,\mathrm{den})$, giving the threshold $\om>1$ and applying to the
odd-denominator shadows; Berth\'e--Holton--Zamboni's formula
$\ice(c_\alpha)=1+\limsup_k q_{k+1}/q_k$ with the elementary floor
$\limsup_k q_{k+1}/q_k\ge\varphi$; and the Bernstein--Lagarias convention
$\PH\circ S\circ\PH^{-1}=T$, from which \eqref{eq:affine} follows. The conjugacy map, the
periodic-value formula and the isometry were re-derived from the parity-vector definition and
checked against the closed forms used here. The height bound, the balanced-numerator bound and the
depth identity were each verified directly at every level the proof uses --- $46$ of $46$ across
four slopes --- rather than by analogy; the slope threshold $\gs$ was checked from both sides on
noble slopes; and controls on a rational target and on the quadratic irrational $\sqrt{17}\in\Zt$
both stayed within $\om\le1+O(1/\log H)$. That session also supplied the witness of
Example~\ref{ex:witness} and the quartic of the footnote in \S\ref{ss:norm}.
\S\ref{sec:measure} was outside its remit.

AI assistance was used for algebra checking, numerical verification, adversarial review, literature
organization, and manuscript drafting. All conceptual development, research direction,
interpretation, and final responsibility for the content belong to the author.

\appendix

\section{Theorem R from Schlickewei's subspace theorem (not new)}\label{app:schlickewei}

Included for completeness only. Theorem R is \cite[Theorem 1.3]{BK18}; the argument below
re-derives it and claims no novelty. It has the mild advantage of being checkable without recourse
to \cite{Ridout58}.

\begin{theoremS}[{$p$-adic subspace theorem; \cite{Schlickewei76,Schlickewei77}, in the formulation
of \cite[Theorem 7.2.2]{BG06}}]
Let $n\ge1$, $K$ a number field, $S$ a finite subset of $M_K$. For each $w\in S$ let
$L_{1,w},\dots,L_{n,w}$ be $n$ linearly independent linear forms in $X_1,\dots,X_n$ with algebraic
coefficients. Then for every $\varepsilon>0$ the set of nonzero $\mathbf{x}\in K^{n}$ with
\[
  \prod_{w\in S}\prod_{i=1}^{n}\frac{|L_{i,w}(\mathbf{x})|_w}{|\mathbf{x}|_w}
  \le H(\mathbf{x})^{-n-\varepsilon}
\]
lies in finitely many proper subspaces of $K^{n}$.
\end{theoremS}

\begin{proposition}\label{prop:Rderiv}
Let $\xi\in\Qt$ be algebraic over $\Q$ and irrational. For every $\varepsilon>0$ there are only
finitely many $r=p/q\in\Q$ in lowest terms with $q>0$ odd and $|\xi-r|_2\le H(r)^{-2-\varepsilon}$.
Equivalently $\om(\xi)\le1$.
\end{proposition}

\begin{proof}
Take $K=\Q$, $n=2$, $S=\{\infty,2\}$, and fix an extension of $|\cdot|_2$ to $\overline{\Q}$
compatible with the given embedding $\Q(\xi)\hookrightarrow\Qt$. Set
\[
  L_{1,2}=X_1-\xi X_2,\quad L_{2,2}=X_2,\qquad L_{1,\infty}=X_1,\quad L_{2,\infty}=X_2 .
\]
At each place the two forms are linearly independent (coefficient determinant $1$) and their
coefficients are algebraic, as Theorem S requires. Let $\mathbf{x}=(p,q)$ with $p,q$ coprime
integers and $q>0$ odd. Then $|\mathbf{x}|_\infty=\max(|p|,q)=H$, while
$|\mathbf{x}|_w=\max(|p|_w,|q|_w)=1$ at every finite place $w$ by coprimality; hence
$H(\mathbf{x})=\prod_w|\mathbf{x}|_w=H$. The double product is
\[
  \text{at }\infty:\ \frac{|p|}{H}\cdot\frac{q}{H}=\frac{|p|q}{H^{2}}\le1,
  \qquad
  \text{at }2:\ |p-\xi q|_2\cdot|q|_2=|p-\xi q|_2=|\xi-p/q|_2
\]
(the last since $q$ is odd, so $|q|_2=1$), so
$\prod_w\prod_i|L_{i,w}(\mathbf{x})|_w/|\mathbf{x}|_w\le|\xi-p/q|_2$. Hence every $r=p/q$ with
$|\xi-r|_2\le H^{-2-\varepsilon}$ satisfies the hypothesis of Theorem S with $n=2$, and all such
$(p,q)$ lie in finitely many proper subspaces of $\Q^{2}$, i.e.\ finitely many lines through the
origin. A line through the origin contains at most one coprime pair $(p,q)$ with $q>0$, hence at
most one rational $p/q$. Finitely many lines, finitely many $r$. Irrationality of $\xi$ gives
$p-\xi q\ne0$, so no form vanishes on the relevant $\mathbf{x}$.
\end{proof}

\section{Exact-arithmetic verification}\label{app:verif}

All computations are in exact integer or rational arithmetic; no floating-point number decides any
inequality. Slopes are pinned by a certified integer interval and every floor or ceiling used to
build a word is checked to agree at both endpoints. Approximation exponents are reported as the
certified interval $[D/\mathrm{bitlen}(H)-1,\;D/(\mathrm{bitlen}(H)-1)-1]$, since
$\mathrm{bitlen}(H)-1\le\log_2 H<\mathrm{bitlen}(H)$. Scripts accompany this paper in
\texttt{verification/}.

\emph{Primitives.} $\PH$ was implemented twice: once from the closed form of
\cite[Proposition 2.1]{DJirr} and once, independently, from the parity-vector definition alone via
the recursion $\PH(v)=2\PH(\sigma v)$ or $(2\PH(\sigma v)-1)/3$ according to $v_0$. The closed
form, the periodic-value formula of Proposition~\ref{prop:shadow}, and the isometry of
Proposition~\ref{prop:iso} were then checked against the independent implementation in $300$, $200$
and $300$ exact random trials respectively; all matched.

\emph{The 2009 series.} Equation \eqref{eq:LS09series} was evaluated in exact $2$-adic arithmetic
and compared with the independently computed $\PH(1c_\alpha)$: agreement was exact modulo $2^{4000}$
at all three slopes tested (golden, $\sqrt2-1$, $\log_3 2$; $16$, $9$, $9$ terms). The $j$-th term
has $\vtwo=q_{j+1}+q_j-1$, which coincides with the depth law at the convergents from above.

\emph{Depth law and height bound across the $\gs$ range (this audit).} The depth identity
$\lcp(c_\gamma,W^{\infty})=q_{k+1}+q_k-2$ and the Step-2 bound \eqref{eq:hbound} were re-checked at
\emph{every} convergent level, both parities, on six slopes spanning the range of
Corollary~\ref{cor:uncond} and crossing $\gs$. Words were generated at $300$ decimal digits, which
for the lengths used exceeds the nearest-integer distance of $j\gamma$ by more than $200$ orders of
magnitude; all heights, $c_W$, denominators and gcds are exact integers.

\begin{center}\footnotesize
\begin{tabular}{@{}lccccc@{}}
\toprule
slope $\gamma$ & $\gamma$ & $\Av$ & $\gamma<\gs$ & depth $=q_{k+1}+q_k-2$ & \eqref{eq:hbound} holds\\
\midrule
$\log_3 2$ (resonance)      & $0.6309298$ & $1$         & yes & all $k$ & all $k$\\
$(\sqrt5-1)/2$              & $0.6180340$ & $1$         & yes & all $k$ & all $k$\\
$\sqrt2-1$                  & $0.4142136$ & $1$         & yes & all $k$ & all $k$\\
$\sqrt3-1=[0;1,2,1,2,\dots]$& $0.7320508$ & $1.1602731$ & yes & all $k$ & all $k$\\
$(\sqrt{21}-3)/2=[0;1,3,1,3,\dots]$ & $0.7912878$ & $1.2541616$ & yes & all $k$ & all $k$\\
$2\sqrt2-2=[0;1,4,1,4,\dots]$       & $0.8284271$ & $1.3130259$ & \emph{no} & all $k$ & all $k$\\
\bottomrule
\end{tabular}
\end{center}

\noindent The last slope lies above $\gs$ and is a control: \eqref{eq:T} still holds there, since
its partial quotients give $\limsup_k q_{k+1}/q_k\ge4$ and hence $\ice\ge5>2\Av=2.626$, with a
measured exponent $1+w$ reaching $4.44$ on the odd levels. This confirms that $\gs$ is sufficient
and not necessary, as \S\ref{ss:aboveg} states. On the slopes with $\Av>1$ the exponent dips below
$2$ at some individual even levels (e.g.\ $1.80$ and $1.95$ for $\sqrt3-1$); this is harmless, the
hypothesis of Theorem~\ref{thm:main} being a limit superior over prefixes and not a condition at
every level.

\emph{The truncation propositions of \S\ref{sec:truncations}.} The determinant identity
$P_NQ_{N+1}-P_{N+1}Q_N=+3^{N}2^{S_N}$ was confirmed in exact integer arithmetic for
$N=1,\dots,120$, with $\vtwo(D_N)/\log_2|D_N|$ rising to $0.49974$ at $N=120$, against the limit
$\frac12$ of Remark~\ref{rem:sunit}. For Proposition~\ref{prop:trunc}, $S_N-\log_2 Q_N$ was found in
$(-1,0]$ at every $N$ tested, while $S_N-\log_2 H(\Xi_N)+\log_2 N$ stayed in $[0.83,1.68]$ for
$N=3,5,10,40,160,640,2560$ --- bounded, as claimed. The crossover $C_N>3^{N}$ first occurs at $N=4$.

\emph{Depth identity and heights.} For the prefixes of $c_\gamma$ at the convergent lengths,
$\lcp(c_\gamma,W^{\infty})=q_{k+1}+q_k-2$ was confirmed at $46$ of $46$ levels across four slopes
(golden $19$, $\sqrt2-1$ $11$, $[0;5,1,5,1,\dots]$ $9$, $\log_3 2$ $7$), the variant with $-1$
failing at $0$ of $46$. The bound $c_W\le3\ell\max(2^{\ell},3^{k})$ and the bound \eqref{eq:hbound}
each held at $46$ of $46$ levels, checked for these prefixes directly rather than by analogy.

\emph{Trend of $\om$.} The certified lower bounds increase to the predicted limit
$\ice/\Av-1$: for the golden slope $1.5654,1.5853,1.5977,1.6028,1.6102,1.6122,1.6151,1.6158,1.6169$
at $k=12,\dots,20$ against the limit $\varphi=1.618034$; for $\sqrt2-1$ to $2.41407$ against
$1+\sqrt2+1=2.414214$; for $[0;5,1,5,1,\dots]$ to $5.853$ against $5.854$; and at $\gamma=\log_3 2$
all seven computed levels exceed $1$, reaching $23.274$.

\emph{Boundary, both sides.} On noble slopes, where $\ice=1+\varphi$ exactly and the margin is
smallest:
\begin{center}
\begin{tabular}{lcccl}
\toprule
$\gamma$ & $\Av$ & predicted limit & measured max & outcome \\
\midrule
$0.61803399$ & $1.000000$ & $1.618034$ & $1.616895$ & $\om>1$ \\
$0.82200179$ & $1.302873$ & $1.009432$ & $1.009305$ & $\om>1$ (inside $\gs$) \\
$0.84889772$ & $1.345503$ & $0.945766$ & $0.945573$ & $\om\le1$ (outside $\gs$) \\
$0.89602865$ & $1.420205$ & $0.843419$ & $0.843273$ & $\om\le1$ \\
\bottomrule
\end{tabular}
\end{center}

\emph{Controls.} On rational targets $\PH(W_0^{\infty})$, with the target itself and all height-$1$
shadows excluded, the worst exponent over all other prefix shadows was $0.25$, $0.60$, $1.00$,
$1.00$ for four test words --- all $\le1$, as they must be, since for fixed rational $\xi$ and
$r\ne\xi$ one has $|\xi-r|_2\ge1/(2H(\xi)H(r))$. On the quadratic irrational $\sqrt{17}\in\Zt$
(computed by bitwise Hensel lifting to $2^{1400}$ and verified), an exhaustive search over all
coprime $p/q$ with $q$ odd and $H\le1200$ gives, band by band, worst exponents
$0.26$, $1.58$, $1.89$, $0.77$, $0.73$, $1.44$, $0.92$, $0.88$, $1.18$, $0.90$ for
$H\in[2^{1},2^{11})$ --- each within the Liouville cap $1+\log_2 9/\log_2 H$ and decaying toward
$1$. The contrast with the previous paragraph is the point of Lemma~\ref{lem:margin}: for an
algebraic target the excess over $1$ \emph{is} the slack and vanishes, while for $\PH(c_\gamma)$ it
converges to a positive constant.

\emph{The parity dichotomy and the witness.} For $1c_\gamma$ the prefix power at the length-$q_n$
prefix equals $1$ at every even $n$, confirmed at every even level of four slopes; an exhaustive
sweep over every prefix length $\ell<900$ found the maximum prefix power attained at $\ell=q_n$ with
$n$ odd and nowhere else. The witness of Example~\ref{ex:witness} was computed as reported there.

\emph{Transfer and sign.} The identities \eqref{eq:affine} were verified exactly modulo $2^{400}$ at
three slopes. The sign check of Remark~\ref{rem:xi} was performed modulo $2^{3000}$ with $1893$
terms.

\bibliographystyle{amsplain}
\bibliography{refs}

\end{document}
